\section*{Chapter \ref{ch:b3:crystalline-solids} --- Crystalline Solids}

\begin{solution}{exo:b3:crystalline-solids:1}
(a) sc: $8\times\tfrac18 = 1$; bcc: $8\times\tfrac18 + 1 = 2$;
fcc: $8\times\tfrac18 + 6\times\tfrac12 = 4$. (b) Diagonal
$a\sqrt3 = 4r$: $r = a\sqrt3/4$. (c) Face diagonal $a\sqrt2 = 4r$:
$r = a\sqrt2/4$. (d) $0.68$ and $0.74$: direction-blind metallic
bonding wants maximal packing, hence the fcc (or the equally dense
hexagonal) structures of copper, aluminium, silver, gold.
\end{solution}

\begin{solution}{exo:b3:crystalline-solids:2}
(a) $a/\sqrt2 = \qty{2.55}{\angstrom}$. (b) 12 --- the
close-packing coordination. (c) \qty{8960}{kg/m^3}
(\cref{ex:b3:crystalline-solids:density}). (d) $n = 4/a^3 =
\qty{8.5e28}{m^{-3}}$, so a cubic micrometre holds
$8.5\times10^{10}$ atoms --- why chip metallurgy is
statistics, not carpentry.
\end{solution}

\begin{solution}{exo:b3:crystalline-solids:3}
(a) $\sin\theta_1 = \lambda/2d = 0.273$: $\theta_1 =
15.8^\circ$. (b) $n \le 2d/\lambda = 3.66$: three orders. (c)
$\sin\theta \le 1$ forces $n\lambda \le 2d$; visible light's
$\lambda \sim \qty{5000}{\angstrom}$ is a thousand times too long
--- no \omterm{def:b3:crystalline-solids:lattice}{crystal} plane spacing can diffract it. (d) $\sin\theta$
falls 1\,\%: $\theta_1 \to 15.7^\circ$. A diffractometer is a fine
thermometer --- and this shift is how thermal expansion is
measured at the atomic scale.
\end{solution}

\begin{solution}{exo:b3:crystalline-solids:4}
(a) $d_{100} = a > d_{110} = a/\sqrt2 > d_{111} = a/\sqrt3$. (b)
Largest $d$, smallest angle: (100). (c) The octahedral (111)
family --- in the diamond structure its sheets pair into strongly
bonded double layers with wide, sparsely bonded gaps between:
the cleaver's plane. (d) Each grain reflects at the same $\theta$
but in a random azimuth: the reflected rays fan into cones of
half-angle $2\theta$, cutting the detector in rings.
\end{solution}

\begin{solution}{exo:b3:crystalline-solids:5}
(a) Interference needs path differences of order $\lambda$;
spacings are ångströms, so $\lambda$ must be too. (b) $E =
hc/\lambda \approx \qty{8}{keV}$. (c) $p = h/\lambda$, $V =
p^2/2m_{\text{e}}e \approx \qty{60}{kV}$ (relativity shaves a few
percent) --- an electron microscope's working voltage. (d)
$\lambda = h/\sqrt{3m_{\text{n}}k_{\text{B}}T} \approx
\qty{1.5}{\angstrom}$: room-temperature neutrons are born
diffraction-ready. They scatter off nuclei, not electron clouds
--- so they see hydrogen clearly and carry a magnetic moment
that maps magnetic order, both nearly invisible to X-rays.
\end{solution}

\begin{solution}{exo:b3:crystalline-solids:6}
(a) Four pairs per cell of volume $a^3$. (b) $N_{\text{A}} =
4M/\rho a^3 = 4\times0.05844/(2165\times\qty{1.794e-28}{})
\approx \qty{6.02e23}{mol^{-1}}$. (c) $N_{\text{A}} \propto
a^{-3}$: a 0.1\,\% error in $a$ is 0.3\,\% in $N_{\text{A}}$. (d)
The method counts atoms by assuming every cell is full and
identical: vacancies, impurities and mosaic boundaries all
miscount --- hence the fanatically perfect silicon spheres of the
kilogram redefinition.
\end{solution}

\begin{solution}{exo:b3:crystalline-solids:7}
(a) Each ion sees $2\sum_{n\ge1}(-1)^{n+1}/n = 2\ln2$ in units of
$e^2/4\pi\varepsilon_0r_0$ (factor 2: both sides), attractive. (b)
$e^2/4\pi\varepsilon_0r_0 = \qty{5.11}{eV}$; $\times1.7476 =
\qty{8.93}{eV}$. (c) Minimising kills $1/9$ of the attraction:
$U = 8.93\times\tfrac89 = \qty{7.94}{eV}$. (d) Within one percent
of \qty{7.9}{eV}: a point-charge \omterm{def:b3:crystalline-solids:lattice}{lattice} plus one stiffness
exponent explains an ionic solid's entire budget --- the quantum
mechanics hides inside $r_0$ and the exponent.
\end{solution}

\begin{solution}{exo:b3:crystalline-solids:8}
(a) $r_0 = 2^{1/6}\times3.40 = \qty{3.82}{\angstrom}$, 1.5\,\%
above the measured \qty{3.76}{\angstrom} (each atom also feels
its twelve neighbours, tightening the well). (b)
$T_{\text{m}} \sim \epsilon/2k_{\text{B}} \approx \qty{60}{K}$:
the right scale for \qty{84}{K}. (c) Helium's \omterm{thm:b3:harmonic-oscillator:spectrum}{zero-point energy}
$\hbar\omega/2$ in so shallow a well rivals the well itself: the
\omterm{def:b3:crystalline-solids:lattice}{crystal} shakes itself apart, and helium stays liquid at $T = 0$
unless squeezed. (d) Depth $\sim\qty{0.01}{eV}$ versus covalent
$\sim\qty{3.6}{eV}$ per bond: three hundredfold in energy is the
whole distance from frost to diamond.
\end{solution}

\begin{solution}{exo:b3:crystalline-solids:9}
(a) $ka \ll 1$: $\omega \approx 2\sqrt{K/m}\,(ka/2) =
a\sqrt{K/m}\,k$. (b) Stretch a cell by $\delta$: stress $\sim
K\delta/a^2$, strain $\delta/a$, so $E \sim K/a$, i.e.\ $K \sim
Ea$. (c) $K \approx \qty{31}{N/m}$, $v = a\sqrt{K/m} \approx
\qty{4400}{m/s}$ --- the measured \qty{4000}{m/s} within 10\,\%,
from a spring guessed off \omterm{thm:b3:continuum-elasticity:hooke}{Young's modulus}. (d) $v = a\sqrt{K/m}$:
stiff bonds up, light atoms up --- diamond maxes both, hence
\qty{18000}{m/s} sound and $\Theta_{\text{D}} \approx
\qty{2200}{K}$.
\end{solution}

\begin{solution}{exo:b3:crystalline-solids:10}
(a) $\dd\omega/\dd k \propto \cos(ka/2) = 0$ at $k = \pi/a$. (b)
$u_n \propto (-1)^n$: every atom in antiphase with both
neighbours --- a standing wave, energy sloshing in place. (c)
With $\lambda = 2a$, waves scattered backwards by successive
atoms differ in path by exactly one wavelength: Bragg's condition
along the chain itself. The \omterm{def:b3:crystalline-solids:lattice}{lattice} reflects its own vibration,
forward and backward waves lock into the standing pattern, and
the travelling wave cannot proceed. (d) $\omega_{\text{max}} = 2\sqrt{K/m} \approx
\qty{3.4e13}{rad/s}$, i.e.\ $\sim\qty{5}{THz}$: the far infrared
--- \omterm{def:b3:crystalline-solids:lattice}{lattice} vibrations and infrared light meet in the same octave,
the fact behind the next exercise.
\end{solution}

\begin{solution}{exo:b3:crystalline-solids:11}
(a) In the optical mode the $+$ and $-$ sublattices move opposite
ways: an oscillating electric dipole, exactly what a light wave
grips. (b) $\omega \approx \sqrt{2K/\mu}$ with $\mu =
\qty{2.3e-26}{kg}$: $\omega \approx \qty{4.7e13}{rad/s}$,
$\lambda = 2\pi c/\omega \approx \qty{40}{\micro m}$ --- the far
infrared (salt's measured reststrahlen band sits at
$\qty{61}{\micro m}$: right octave from a two-spring model). (c)
At that band the \omterm{def:b3:crystalline-solids:lattice}{crystal}'s own resonance absorbs and re-reflects
the wave: transparent salt turns mirror-opaque. (d) One atom per
cell means no second sublattice to beat against: a single branch,
no gap.
\end{solution}

\begin{solution}{exo:b3:crystalline-solids:12}
(a) $N$ atoms $\times$ 3 displacement directions = $3N$
oscillators, so the linear spectrum is cut off once it has
counted $3N$ modes: that defines $k_{\text{D}}$. (b)
$k_{\text{D}} = (6\pi^2n)^{1/3} = \qty{1.7e10}{m^{-1}}$,
$\Theta_{\text{D}} = \hbar vk_{\text{D}}/k_{\text{B}} \approx
\qty{350}{K}$ --- the calorimetric \qty{343}{K}, from a sound
speed and a density. (c) The straight line misses the flattening
at the zone edge: Debye over-counts high frequencies, so the fit
strains at intermediate temperatures; the $T^3$ law (long waves
only) is safe. (d) Both say the spectrum ends when the wavelength
reaches the atomic spacing --- one mode per atom, dressed either
as a cutoff wavevector or a cutoff temperature.
\end{solution}

\begin{solution}{pb:b3:crystalline-solids:1}
\textbf{1.} $2d\sin\theta = n\lambda$. Monochromatic: one
$\lambda$ makes each ring angle map to one spacing. Powder: among
random grains, some are always oriented to reflect --- every
family speaks at once.
\textbf{2.} $\theta = 19.05^\circ, 22.15^\circ, 32.2^\circ,
38.75^\circ$; $\sin^2\theta = 0.107, 0.142, 0.284, 0.392$.
\textbf{3.} Insert $d = a/\sqrt{h^2+k^2+l^2}$ into Bragg
($n = 1$): $\sin^2\theta = (\lambda^2/4a^2)(h^2+k^2+l^2)$.
\textbf{4.} Ratios $1 : 1.33 : 2.67 : 3.68$; times 3:
$3.0 : 4.0 : 8.0 : 11.0$.
\textbf{5.} All-odd $(111)$ and all-even $(200)$, $(220)$,
$(311)$ give exactly 3, 4, 8, 11; the absent 1, 2, 5, 6, 7 rule
out simple cubic, whose pattern would show them.
\textbf{6.} Bcc's allowed even sums give ratios
$1:2:3:4$ --- a different rhythm. Ratios survive not knowing
$a$, wavelength drift, and sample misalignment: parity is
geometry, not calibration.
\textbf{7.} Few planes make a blunt interference maximum, like a
grating with five slits: ring width $\sim$ inverse crystallite
size (the Scherrer relation) --- sharp rings certify well-grown
grains.
\textbf{8.} $a = \lambda\sqrt{h^2+k^2+l^2}/2\sin\theta = 4.088,
4.086, 4.089, 4.082\,\unit{\angstrom}$.
\textbf{9.} $a \approx \qty{4.09}{\angstrom}$.
\textbf{10.} $\rho = 4M/N_{\text{A}}a^3$.
\textbf{11.} $4/N_{\text{A}}a^3 = \qty{9.74e4}{mol/m^3}$:
silver $\to \qty{10.5}{g/cm^3}$ (matches silver's handbook
value), gold $\to 19.2$ (matches gold's!), platinum $\to 19.0$
(contradicts platinum's 21.4 --- its true $a$ is
\qty{3.92}{\angstrom}). Platinum is out; silver and gold both
survive, because their \omterm{def:b3:crystalline-solids:lattice}{lattice constants} coincide almost exactly.
\textbf{12.} A mass-based measurement: density (or simply molar
mass). Silver and gold share $a$ to 0.2\,\%, but their atomic
masses differ by a factor 1.8 --- no \omterm{def:b3:crystalline-solids:lattice}{lattice} coincidence can
bridge that.
\textbf{13.} Archimedes on the ingot: near \qty{10.5}{g/cm^3}
$\to$ silver, consistent with the corrosion (silver tarnishes;
gold would have come up gleaming). Microscopic cell mass and
macroscopic weighing agree --- the same
$\rho = 4M/N_{\text{A}}a^3$ read in both directions.
\textbf{14.} The atoms rattle about \emph{unmoved average
positions}: the mean \omterm{def:b3:crystalline-solids:lattice}{lattice}, which fixes ring angles, is intact;
random displacements only redistribute intensity.
\textbf{15.} Intensity falls smoothly with $T$ (a
$\exp(-q^2\langle x^2\rangle)$-type Debye--Waller factor), the
lost light reappearing as diffuse background between rings.
\textbf{16.} $K \approx Ea = \qty{8.3e10}{}\times
\qty{2.89e-10}{} \approx \qty{24}{N/m}$.
\textbf{17.} $v = a\sqrt{K/m} \approx \qty{3300}{m/s}$ ---
the measured \qty{2700}{m/s} to 20\,\%.
\textbf{18.} $\langle x^2\rangle = 3k_{\text{B}}T/K$: rms
$\approx\qty{0.23}{\angstrom}$, about 8\,\% of the bond length
already at room temperature.
\textbf{19.} The 10\,\% mark extrapolates to $\sim\qty{500}{K}$
against the real \qty{1235}{K}: right order, factor two-ish out
--- our single-spring $K$ is too soft, and Lindemann is a scaling
rule, not a law. The lesson stands: melting is when thermal
rattle rivals the \omterm{def:b3:crystalline-solids:lattice}{lattice} itself.
\textbf{20.} $d_{110} = a/\sqrt2 = \qty{2.04}{\angstrom}$:
$\sin\theta = 0.378$, $2\theta \approx 44.5^\circ$.
\textbf{21.} Coupling strength: electrons (charged) scatter in
nanometres --- surfaces and foils; X-rays in micrometres ---
bulk powder; neutrons in centimetres --- whole engine parts.
\textbf{22.} The rings collapse into one or two broad halos:
the liquid keeps short-range order (preferred neighbour
distance) but no long-range register --- nothing periodic left to
interfere sharply.
\textbf{23.} DNA --- Photo 51. Diffraction reads repeat
distances far below any light microscope's reach; the X pattern
betrayed a helix, the \qty{3.4}{\angstrom} smear its stacked
rungs.
\textbf{24.} That the structure is deterministically ordered
(quasiperiodic --- ordered without repeating): sharp spots need
long-range phase coherence, not periodicity --- the discovery
that widened crystallography's own definition of a \omterm{def:b3:crystalline-solids:lattice}{crystal}.
\textbf{25.} Ratios $3:4:8:11 \to$ fcc; $a = \qty{4.09}{\angstrom}$
$+$ density $\to$ silver, not gold or platinum; and since heating
fades rings and melting erases them, X-ray the ingot cold ---
the museum's silver, certified by interference.
\end{solution}
